\begin{solapartat}
\figura[0.25]{sol_fig1.png}

\punts{0,1} La força aplicada pel camp magnètic és:
\[ \vec{F}_M = q\vec{v}\times\vec{B} \]

\punts{0,2} Aplicant la regla de la mà dreta i tenint en compte que els ions estan carregats negativament, podem veure que la força magnètica està dirigida verticalment cap amunt.

\punts{0,2} Perquè l'ió no es desviï, la suma de forces ha de ser zero, $\vec{F}_E + \vec{F}_M = \vec{0}$, per tant la força aplicada pel camp elèctric ha de ser vertical i apuntant cap avall.

\punts{0,1} La força aplicada pel camp elèctric és:
\[ \vec{F}_E = q\vec{E} \]

\punts{0,1} Com que la càrrega és negativa, el camp elèctric ha de ser vertical i apuntant cap amunt.

\punts{0,2} Perquè no es desviïn els mòduls de les forces elèctrica i magnètica han de ser iguals:
\[ qE = qvB \;\Rightarrow\; E = vB = 10^{5}\ \frac{\mathrm{V}}{\mathrm{m}} \]

\punts{0,1} El camp creat per dues plaques paral·leles és homogeni i la relació entre diferència de potencial i el camp és:
\[ \Delta V = Ed \]
On $d$ és la distància entre plaques, per tant la diferència de potencial que cal aplicar és:
\[ \Delta V = Ed = 10^{5}\cdot 0,015 = 1500\ \mathrm{V} \]

\punts{0,1} El camp elèctric apunta en la direcció en la qual el potencial disminueix; per tant, cal connectar la placa inferior al terminal de potencial alt i la placa superior al terminal de potencial baix:

\noindent\begin{minipage}[c]{0.62\linewidth}
\figura[0.95]{sol_fig2.png}
\end{minipage}\hfill
\begin{minipage}[c]{0.3\linewidth}
\figura[0.95]{sol_fig3.png}
\end{minipage}

\punts{0,15} Aquest selector de velocitats configurat així també funciona per càrregues positives atès que tant la força elèctrica com la magnètica canvien el sentit, per tant, la seva suma vectorial segueix sent zero.
\end{solapartat}

\begin{solapartat}
\punts{0,2} representació

\figura[0.75]{sol_fig4.png}

\punts{0,3} A més, la força magnètica és perpendicular a $\vec{v}$, per tant, el resultat serà que aplicarem una força normal constant; és a dir, l'ió descriurà un moviment circular uniforme (mòdul de la velocitat constant). El sentit de la força ve determinat per la regla de la mà dreta.

\punts{0,5} La força total és la força aplicada pel camp magnètic:
\[ \vec{F}_M = -e\vec{v}\times\vec{B} \]
El radi de la trajectòria es pot obtenir a partir de la segona llei de Newton:

$\vec{F} = m\vec{a}$

i tenint en compte que l'acceleració centrípeta s'expressa com
\[ a = \frac{v^2}{r} \]
I, finalment, com la velocitat i el camp magnètic són perpendiculars:
\[ evB = m\,\frac{v^2}{r} \;\Rightarrow\; r = \frac{mv}{eB} \]

\punts{0,25} I la distància serà dues vegades el radi, per tant, per l'ió $^{20}Ne^{-}$ tenim:
\[ d = 2\,\frac{mv}{eB} = 2\,\frac{3,32\times 10^{-26}\times 2\times 10^{5}}{1,602\times 10^{-19}\times 0,2} = 0,414\ \mathrm{m} \]
\end{solapartat}
